% Pursuing Stacks (cotes 134-2 to 134-8): archive cote 134-2, site page 44 (PDF page 45).
% appendix (letters to Breen). Draft compared: ps1.tex:1212-1216.
% AI-transcribed by Claude Opus 5.5 (claude-opus-5-5, adaptive thinking, effort high), Message Batches API, batch pilot (msgbatch_0189Nhf1xB615YWnztubtLfc), 2026-11-12; not yet checked by eye.
\scanpage{134-2/44}
\mspage{29}

morphisme sur les $\hand{\pi}_0$ et sur les $\hand{\pi}_i$ en tout point (1\hand{$\leq$}i\hand{$\leq$}n). On serait plus
heureux de pouvoir dire à la place ``et \inserted{de plus} un homomorphisme surjectif pour
i=\over{N}{n}+1'', car c'est, il me semble, cela qu'il faudrait pour \inserted{espérer} pouvoir conclure
que la catégorie localisée de la catégorie des ensembles semi-simpliciaux,
obtenue en inversant les flèches\hand{\textbackslash}``qui induisent des isomorphismes sur les
$\hand{\pi}_i$ pour 0\hand{$\leq$}i\hand{$\leq$}n'' (ou encore,``en négligeant''les ensembles semi-simpliciaux
n-connexes), est équivalente à la catégorie localisée de la catégorie des
n- catégories, où on rend inversibles les n-équivalences ? Quoi qu'il en
soit, ces petites bavures devraient disparaitre lorsqu'on ``stabilise'' en
faisant augmenter n. A ce propos, on voit que le foncteur ``tron\over{\illegible{1}}{ca}ture en
dim. n'' de la théorie homotopique (consistant à tuer les groupes d'homotopi
% NOTE: the end of this line is cut off at the right edge of the scan ("homotopi...")
à partir de la dimension n\unsure{\hand{+1)}}.s'interprète dans le langage des n-catégories
par l'opération faisant passer d'une N-catégorie (N\hand{$>$}n) à une n-catégorie,
en conservant \xout{\illegible{2}} tels quels les i-objets (0\hand{$\leq$}i\hand{$\leq$}n-1) et leur
composition (1\hand{$\leq$}i\hand{$\leq$}n-1), et en remplaçant les n-objets par les classes \xout{d'}
\xout{\unsure{équivalence}} de n-objets ``à équivalence près'', avec la composition obtenue
par passage au quotient. De même, le foncteur d'inclusion évident en
théorie homotopique, consistant à regarder un ensemble semi-simplicial
``où on a négligé les $\hand{\pi}_i$ pour i\hand{$\geq$}n+1'' comme un ensemble semi-simplicial \xout{qui}
(dans la catégorie homotopique) qui se trouve avoir des $\hand{\pi}_i$ nuls pour
i\hand{$\geq$}n+1, se traduit par le foncteur allant des n-catégories vers les N-%
catégories, obtenue en ajoutant à une n-catégorie des i-flèches (n+1 \hand{$\leq$}i\hand{$\leq$}N)
identiques exclusivement. (Ainsi, un ensemble est regardé comme une catégo-%
rie ``discrète'', une catégorie comme une 2-catégorie où les \uline{Hom}(A,B), A et B
des 0-objets, sont des catégories discrètes, etc ...).

Bien entendu, rien n'empêche de considérer aussi la notion de $\hand{\infty}$-catégo-%
% NOTE: line end cut at the right edge; a hyphen after "catégo" may be hidden
rie, à laquelle celle de n-catégorie est comme la notion d'ensemble semi-\xout{x}
simplicial tronqué à celle d'ensemble semi-simplicial. Sauf erreur, la
localisée de la catégorie des $\hand{\infty}$-catégories, pour les\xout{x} flèches de $\hand{\infty}$-équiva-%
% NOTE: line end cut at the right edge; a hyphen after "équiva" may be hidden
lence, est \inserted{équivalente à} ``la catégorie homotopique''\hand{,} localisée de la catégorie des ensemble\hand{s}
semi-simpliciaux, ou du moins une sorte de complétée de celle-là. Dans
cette optique, le tapis consistant à interpréter une $\hand{\infty}$-catégorie de Picard
stricte (i.e. quelque chose qui ressemble à un groupe abélien de la catégo-%
rie des $\hand{\infty}$-catégories) comme donnée (à $\hand{\infty}$-équivalence près) par un \xout{objet}
complexe de chaines regardé comme un objet d'une catégorie dérivée, est à
relier au tapis de Dold-Puppe, interprétant ces derniers comme des groupes
abéliens semi-simpliciaux.

Pour se donner confiance dans ce yoga\inserted{général}, on peut essayer d'interpéter
en termes de n-catégories \xout{des} ou $\hand{\infty}$-catégories des constructions familières

\margin{44}
